\documentclass[12pt,a4paper]{article}
\usepackage{fontspec}
\usepackage{geometry}
\usepackage{enumitem}
\usepackage{fancyhdr}
\usepackage{lastpage}
\usepackage{hyperref}

% Configure IBM Plex Sans Devanagari for mixed Hindi/English content
% IBM Plex Sans Devanagari includes both Latin and Devanagari scripts
\setmainfont{IBMPlexSansDevanagari}[
  Path=/home/asi0/asi0-repos/bitcoin-educational-content/scripts/quizz_to_pdfs/fonts/TrueType/IBM-Plex-Sans-Devanagari/,
  Extension=.ttf,
  UprightFont=*-Regular,
  BoldFont=*-Bold,
  Scale=1.0
]

\geometry{margin=2.5cm}
\pagestyle{fancy}
\fancyhf{}
\rhead{\sffamily Page \thepage\ of \pageref{LastPage}}
\lhead{\sffamily BIZ101 - Quiz (hi)}
\renewcommand{\headrulewidth}{0.4pt}

\title{Quiz: BIZ101}
\author{Language: hi}
\date{\today}

\begin{document}

\maketitle
\thispagestyle{fancy}

\section*{Instructions}
Answer all questions by selecting the best option (A, B, C, or D). The answer key is provided at the end of this document.

\bigskip


\subsection*{Question 1}
पारंपरिक मुद्रा प्रणालियों पर बने पेमेंट रेल्स क्या हैं?

\begin{enumerate}[label=\Alph*., leftmargin=*]
\item कार्ड कंपनियों द्वारा भुगतान किया गया जो ग्राहकों का डेटा दोबारा बेचती हैं
\item ग्राहक द्वारा दुकानदार की मांगी गई कीमत के ऊपर अदा किया गया
\item मूल रूप से भुगतान प्रोसेसरों के शुल्क के माध्यम से व्यापारियों द्वारा भुगतान किया गया
\item मूल रूप से मुफ़्त क्योंकि ग्राहक व्यापारी द्वारा दिखाए गए दाम का भुगतान करता है
\end{enumerate}

\bigskip


\subsection*{Question 2}
कमी को एक अच्छी करेंसी की महत्वपूर्ण खूबी क्यों माना जाता है?

\begin{enumerate}[label=\Alph*., leftmargin=*]
\item क्योंकि इससे हर कोई आसानी से अपनी खुद की करेंसी बना सकता है।
\item क्योंकि इससे करेंसी को पहचानना मुश्किल हो जाता है।
\item क्योंकि यह लंबी दूरी के व्यापार में मुद्रा के उपयोग को हतोत्साहित करता है।
\item क्योंकि सीमित Supply अपनी वैल्यू को ज़्यादा उत्पादन से कम होने से बचाता है।
\end{enumerate}

\bigskip


\subsection*{Question 3}
सिक्कों से नोटों की तरफ बदलाव ने करेंसी नेटवर्क को कैसे बेहतर बनाया?

\begin{enumerate}[label=\Alph*., leftmargin=*]
\item इसने सरकारों को पैसे Supply का बिल्कुल भी प्रबंधन करने से बचा लिया।
\item इसने तुरंत महंगाई को रोक दिया और अर्थव्यवस्था को स्थिर कर दिया।
\item इसने पैसे को और ज्यादा पोर्टेबल और स्केलेबल बना दिया, जिससे लंबी दूरी के व्यापार और जटिल वित्तीय प्रणालियों को आसानी हो गई।
\item इसने किसी भी अंतर्निहित संपत्ति या ट्रस्ट की जरूरत को खत्म कर दिया।
\end{enumerate}

\bigskip


\subsection*{Question 4}
आधुनिक केंद्रीकृत मुद्रा प्रणालियों का एक नुकसान क्या है?

\begin{enumerate}[label=\Alph*., leftmargin=*]
\item वे सभी बिचौलियों को हटा देते हैं और भुगतान को मिलीसेकंड में पूरा कर देते हैं।
\item वे यह सुनिश्चित करते हैं कि पैसे का जारी होना पूरी तरह से विकेंद्रीकृत हो और मुद्रास्फीति से सुरक्षित हो।
\item वे अक्सर पैसे Supply के प्रबंधन में पारदर्शिता और ऑडिट करने की क्षमता की कमी रखते हैं।
\item वे हर मुद्रा इकाई के लिए एक स्थिर, अपरिवर्तनीय मूल्य प्रदान करते हैं।
\end{enumerate}

\bigskip


\subsection*{Question 5}
आज के रिटेल पॉइंट ऑफ सेल (POS) पर कार्ड-आधारित भुगतान प्रक्रिया की जटिलता क्या है?

\begin{enumerate}[label=\Alph*., leftmargin=*]
\item यह पूरी तरह से नकदी पर निर्भर करता है और बीच में कोई दलाल नहीं होता।
\item इसमें सिर्फ एक ही पल का कदम चाहिए और पैसे तुरंत सेटल हो जाते हैं।
\item इसमें कई बिचौलिए शामिल होते हैं, कई चरण होते हैं, और पूरी तरह से फंड सेटल होने में दिनों का समय लग सकता है।
\item यह सुनिश्चित करता है कि जैसे ही कार्ड स्वाइप किया जाता है, व्यापारी को तुरंत और अपरिवर्तनीय रूप से धनराशि मिल जाती है।
\end{enumerate}

\bigskip


\subsection*{Question 6}
Bitcoin की करेंसी नेटवर्क के रूप में एक खास बात यह है कि यह बिना किसी सेंट्रल अथॉरिटी के काम करता है - यानी यह डिसेंट्रलाइज्ड है। इसमें सारे ट्रांजैक्शन पीयर-टू-पीयर होते हैं और किसी बैंक या सरकार का इसमें दखल नहीं होता।

\begin{enumerate}[label=\Alph*., leftmargin=*]
\item यह किसी भी यूजर के अनुरोध पर करेंसी के नए यूनिट जारी करता है।
\item इसे हर लेन-देन को सत्यापित और निपटाने के लिए बैंकों की जरूरत होती है।
\item यह सिर्फ तभी काम कर सकता है जब सरकारी एजेंसी द्वारा अनुमति दी जाए।
\item यह नकदी की तरह है और लेन-देन के लिए किसी केंद्रीय प्राधिकरण या बिचौलिए की जरूरत नहीं होती।
\end{enumerate}

\bigskip


\subsection*{Question 7}
Bitcoin के 21 मिलियन के इश्यू कैप का क्या महत्व है?

\begin{enumerate}[label=\Alph*., leftmargin=*]
\item यह कैप पहुँचने के बाद किसी को भी लेन-देन भेजने से रोकता है।
\item यह Bitcoin के दाम में होने वाले उतार-चढ़ाव को पूरी तरह खत्म कर देता है।
\item यह एक अपस्फीतिकारी (deflationary) संपत्ति बनाता है, जो मनमाने ढंग से मुद्रास्फीति (inflation) के खिलाफ अपने मूल्य की रक्षा करता है।
\item यह सुनिश्चित करता है कि सरकारें जितने चाहें उतने बिटकॉइन छाप सकती हैं।
\end{enumerate}

\bigskip


\subsection*{Question 8}
निम्नलिखित में से कौन सा विकल्प Bitcoin लेनदेन की फाइनैलिटी (अंतिमता) को सबसे बेहतर तरीके से बताता है?

\begin{enumerate}[label=\Alph*., leftmargin=*]
\item उपयोगकर्ताओं को लेन-देन पूरा करने के लिए कई सरकारी संस्थाओं से मंजूरी लेनी होगी।
\item बैंक लेन-देन को वापस ले सकते हैं अगर कोई विवाद हो जाए।
\item Blockchain पर एक बार रिकॉर्ड हो जाने के बाद, लेन-देन को बदला नहीं जा सकता।
\item कोई भी भागीदार पिछले लेन-देन को मनमर्जी से बदल सकता है, अगर उनके पास एक खास कोड हो।
\end{enumerate}

\bigskip


\subsection*{Question 9}
Bitcoin का कारोबारों के लिए नेटवर्क और संपत्ति के रूप में क्या भूमिका है, यह कौन सा कथन सबसे अच्छे से बताता है?

\begin{enumerate}[label=\Alph*., leftmargin=*]
\item इसे चलाने के लिए एक ही देश की मौद्रिक नीतियों से जोड़ा जाना चाहिए।
\item यह सिर्फ एक अटकलबाज़ी वाला ट्रेडिंग टूल है, जिसका कोई असली फायदा नहीं है।
\item यह सिर्फ एक लोकल गिफ्ट कार्ड सिस्टम की तरह काम कर सकता है।
\item यह एक वैश्विक, विकेंद्रीकृत प्रणाली में मूल्य के भंडार और Exchange के माध्यम दोनों के रूप में काम कर सकता है।
\end{enumerate}

\bigskip


\subsection*{Question 10}
Bitcoin के ऊपर बनाए गए Lightning Network का मुख्य उद्देश्य क्या है?

\begin{enumerate}[label=\Alph*., leftmargin=*]
\item लेन-देन को धीमा करना ताकि सरकारें उन्हें पूरी तरह से ट्रैक कर सकें।
\item Bitcoin के Blockchain को पूरी तरह से बदलने और सभी लेन-देन को स्थानीय रूप से स्टोर करने के लिए।
\item एक ही संस्था द्वारा नियंत्रित एक अनुमति-आधारित नेटवर्क प्रदान करना।
\item तेज़, लगभग तुरंत भुगतान सक्षम करने के लिए और On-Chain के साथ न्यूनतम संपर्क के साथ।
\end{enumerate}

\bigskip


\subsection*{Question 11}
Lightning Network कैसे Bitcoin ट्रांजैक्शन के फीस को कम करने में मदद करता है?

\begin{enumerate}[label=\Alph*., leftmargin=*]
\item भुगतान से सारी क्रिप्टोग्राफिक पुष्टि हटाकर।
\item ज्यादातर लेन-देन off-chain पर करके और सिर्फ चैनल सेटलमेंट लेन-देन को मुख्य Blockchain पर पब्लिश करके।
\item एक बड़े पेमेंट हब में नियंत्रण को केंद्रीकृत करके।
\item व्यापारियों को सभी नेटवर्क लागत पहले से चुकाने के लिए बाध्य करके।
\end{enumerate}

\bigskip


\subsection*{Question 12}
लाइटनिंग मर्चेंट्स को पेमेंट फाइनैलिटी के मामले में क्या फायदा देता है?

\begin{enumerate}[label=\Alph*., leftmargin=*]
\item यह लेन-देन को पूरी तरह से 90 दिनों तक वापस लेने योग्य बना देता है।
\item लाइटनिंग पर भुगतान तुरंत पूरा हो जाता है, जिससे रिवर्सल या धोखाधड़ी का खतरा कम हो जाता है।
\item यह ग्राहकों को कभी भी अपना भुगतान वापस लेने की अनुमति देता है।
\item इसमें पैसे खर्च करने से पहले 24 घंटे का इंतज़ार करना पड़ता है।
\end{enumerate}

\bigskip


\subsection*{Question 13}
Lightning Network को "भुगतान के लिए ईमेल" की तरह क्यों कहा जा सकता है?

\begin{enumerate}[label=\Alph*., leftmargin=*]
\item क्योंकि यह पेमेंट्स को पारंपरिक ईमेल प्रदाताओं के जरिए भेजता है।
\item क्योंकि यह बिना किसी रुकावट, वैश्विक और बिना अनुमति के मूल्य Exchange को संदेश भेजने जितनी आसानी से ट्रांसफर करने में सक्षम बनाता है।
\item क्योंकि यह सिर्फ तभी काम करता है जब यह नियमित ईमेल संदेशों से जुड़ा होता है।
\item क्योंकि यह यूज़र के बैलेंस को ईमेल के बॉडी में स्टोर करता है।
\end{enumerate}

\bigskip


\subsection*{Question 14}
ऑस्ट्रियन इकोनॉमिक्स के मुताबिक, मुनाफ़ा किसी कंपनी की सेहत का अहम संकेतक क्यों माना जाता है?

\begin{enumerate}[label=\Alph*., leftmargin=*]
\item क्योंकि मुनाफ़ा दिखाता है कि व्यापार ग्राहकों की ज़रूरतों को कुशलता और टिकाऊ तरीके से पूरा कर रहा है।
\item क्योंकि मुनाफ़ा यह सुनिश्चित करता है कि कंपनी का कर्मचारी दल उद्योग में सबसे बड़ा हो।
\item क्योंकि मुनाफ़ा बाज़ार की अंधाधुंध घटनाओं को दिखाता है जो ग्राहकों की मांग से कोई वास्ता नहीं रखतीं।
\item क्योंकि मुनाफ़ा सिर्फ़ इस बात का संकेत है कि कंपनी को सरकारी सब्सिडी मिली है।
\end{enumerate}

\bigskip


\subsection*{Question 15}
क्योंकि भविष्य अनिश्चित है, इसलिए किसी भी कंपनी के लिए पूंजी को सुरक्षित रखना बेहद ज़रूरी है। अगर कंपनी के पास पर्याप्त पूंजी होगी, तो वह मुश्किल समय में भी अपने कामकाज को चालू रख सकेगी, नए मौकों का फायदा उठा सकेगी और अचानक आने वाले झटकों को झेल सकेगी। बिना पूंजी के, कंपनी का टिक पाना मुश्किल हो जाता है, खासकर जब बाज़ार में उतार-चढ़ाव हो या आर्थिक मंदी आ जाए। इसलिए, पूंजी बचाकर रखना कंपनी की सेहत के लिए बिल्कुल ज़रूरी है।

\begin{enumerate}[label=\Alph*., leftmargin=*]
\item पूंजी को सुरक्षित रखने से यह गारंटी मिलती है कि कोई भी प्रतियोगी बाजार में कभी प्रवेश नहीं कर सकता।
\item कैपिटल बचाना तभी जरूरी है जब कंपनी कभी बढ़ने की प्लानिंग ही न करे।
\item पूंजी भंडार बनाए रखने से कंपनी को ढलने, विकास में निवेश करने और आर्थिक मंदी का सामना करने में मदद मिलती है।
\item पूंजी को सुरक्षित रखने से सभी नियामक परिवर्तनों से स्थायी सुरक्षा मिलती है।
\end{enumerate}

\bigskip


\subsection*{Question 16}
मछुआरे की कहानी पूंजी की प्रकृति के बारे में क्या दिखाती है?

\begin{enumerate}[label=\Alph*., leftmargin=*]
\item सरकारी कर्ज़े से पूंजी तुरंत बनानी चाहिए।
\item पूंजी बचाए गए संसाधनों को दर्शाती है जो भविष्य की उत्पादकता बढ़ाने के लिए औजार बनाने में मदद करते हैं।
\item बड़ी कंपनियों के लिए ही पूंजी मायने रखती है, छोटे उद्यमियों के लिए नहीं।
\item पूंजी सिर्फ उपभोक्ता सामानों का एक ढेर है जो तुरंत खपत के लिए इंतज़ार कर रहा है।
\end{enumerate}

\bigskip


\subsection*{Question 17}
कर्ज़-आधारित पैसे से बनाया गया कृत्रिम पूंजी और ऑस्ट्रियाई स्टाइल के पूंजी में क्या फर्क है?

\begin{enumerate}[label=\Alph*., leftmargin=*]
\item ऑस्ट्रियन-स्टाइल की पूंजी तब बनती है जब जरूरत पड़ने पर ज्यादा पैसा छापा जाता है।
\item ऋण-आधारित पूंजी को कभी भी चुकाने की ज़रूरत नहीं होती और यह असली बचत के बराबर होती है।
\item ऑस्ट्रियन-स्टाइल पूंजी असली बचत और त्याग से आती है, जबकि कर्ज़-आधारित पूंजी क्रेडिट से तुरंत बन जाती है और आर्थिक संकेतों को गड़बड़ा सकती है।
\item कोई फर्क नहीं है; दोनों तरह की पूंजी प्रकृति और प्रभाव में एक जैसी हैं।
\end{enumerate}

\bigskip


\subsection*{Question 18}
एक कंपनी अपने ट्रेजरी का कुछ हिस्सा Bitcoin में क्यों रखना चाहेगी?

\begin{enumerate}[label=\Alph*., leftmargin=*]
\item तुरंत सारा जोखिम खत्म करने और स्थिर रिटर्न्स की गारंटी देने के लिए।
\item सरकारी नियमों का पालन करने के लिए जीडब्ल्यू-1 होल्डिंग्स अनिवार्य हैं।
\item एक संभावित रूप से मूल्यवर्धित संपत्ति के रूप में जो मूल्य संचय कर सकती है और पारंपरिक वित्तीय मध्यस्थों पर निर्भर किए बिना तरलता प्रदान कर सकती है।
\item यह सुनिश्चित करने के लिए कि वह अपनी संपत्ति को कभी भी पारंपरिक मुद्राओं में वापस नहीं बदल सके।
\end{enumerate}

\bigskip


\subsection*{Question 19}
एक कारण क्यों संस्थागत निवेशक और बड़ी कंपनियां Bitcoin में दिलचस्पी लेने लगी हैं?

\begin{enumerate}[label=\Alph*., leftmargin=*]
\item Bitcoin को बड़े वित्तीय खिलाड़ियों जैसे ब्लैकरॉक ने अपनाया है और स्पॉट Bitcoin ETFs की सफलता ने इसे वैल्यू स्टोर करने वाली संपत्ति के रूप में मान्यता दिला दी है।
\item इसमें डिजिटल एसेट्स या वित्तीय नियमों की कोई समझ नहीं चाहिए।
\item जीडब्ल्यू-2 की कीमत सरकारी एजेंसियों द्वारा तय और गारंटीशुदा है।
\item जीडब्ल्यू-3 लेनदेन केंद्रीय बैंकों द्वारा पूरी तरह से बीमाकृत हैं।
\end{enumerate}

\bigskip


\subsection*{Question 20}
वोलेटिलिटी को Bitcoin के लिए एक नेचुरल एस्पेक्ट क्यों माना जाता है, जो एक इमर्जिंग एसेट है?

\begin{enumerate}[label=\Alph*., leftmargin=*]
\item क्योंकि Bitcoin अभी भी अपेक्षाकृत नया है, सट्टेबाजों को आकर्षित करता है और बाजार की खबरों पर प्रतिक्रिया देता है, जिससे कीमतों में उतार-चढ़ाव होता है, लेकिन लंबे समय में यह स्थिर हो जाता है।
\item क्योंकि यह केंद्रीय रूप से नियंत्रित है और कीमतें एक ही अधिकारी द्वारा तय की जाती हैं।
\item क्योंकि Bitcoin में कोई भी यूज़र लंबे समय तक इसे होल्ड नहीं कर रहा।
\item क्योंकि दुकानदारों को हर घंटे कीमतें अप्रत्याशित तरीके से तय करनी पड़ती हैं।
\end{enumerate}

\bigskip


\subsection*{Question 21}
एक डेट-बेस्ड मोनेटरी सिस्टम के कॉन्टेक्स्ट में "हिडन इन्फ्लेशन" क्या होता है?

\begin{enumerate}[label=\Alph*., leftmargin=*]
\item यह Supply में पैसे की बढ़ोतरी (लगभग 7% सालाना) है जो आपके कुल मूल्य के हिस्से को कम कर देती है, अगर आप उसी हिसाब से निवेश या संपत्ति नहीं बढ़ाते।
\item यह उन मुद्रास्फीति के आंकड़ों की बात करता है जो केंद्रीय बैंकों द्वारा प्रकाशित नहीं किए जाते।
\item इसका मतलब है कि कीमतें कभी नहीं बढ़ती, चाहे पैसा Supply बढ़ जाए।
\item यह फिक्स्ड इंटरेस्ट रेट्स के लिए एक आधिकारिक आर्थिक शब्द है।
\end{enumerate}

\bigskip


\subsection*{Question 22}
एक कंपनी Bitcoin को पारंपरिक "सुरक्षित" संपत्तियों जैसे रियल एस्टेट पर क्यों चुन सकती है?

\begin{enumerate}[label=\Alph*., leftmargin=*]
\item जीडब्ल्यू-1 असली महँगाई से आगे निकल सकता है और रियल एस्टेट की अशक्तता और रखरखाव लागत से जुड़ा नहीं है।
\item जीडब्ल्यू-3 सरकारी गारंटी प्रदान करता है जो रियल एस्टेट नहीं देता।
\item रियल एस्टेट हमेशा Bitcoin से सस्ता होता है और उसमें मेंटेनेंस की भी जरूरत नहीं पड़ती।
\item रियल एस्टेट Bitcoin की तुलना में हर मार्केट कंडीशन में तेजी से प्राइस बढ़ाता है।
\end{enumerate}

\bigskip


\subsection*{Question 23}
कंपनी के ट्रेजरी का एक हिस्सा Bitcoin में आवंटित करते समय आमतौर पर क्या सलाह दी जाती है?

\begin{enumerate}[label=\Alph*., leftmargin=*]
\item सारे ऑपरेशनल कैश को तुरंत खर्च कर दो, लिक्विडिटी की जरूरतों की परवाह किए बिना।
\item जब तुम्हें जरूरत हो तभी निवेश करो, वो भी तब जब तुरंत छोटे समय के लिए कर्ज चाहिए खर्चे पूरे करने के लिए।
\item केवल उस अतिरिक्त नकदी को निवेश करें जिसकी आपको कई सालों तक ज़रूरत नहीं होगी और कम से कम चार साल का समय सीमा बनाए रखें।
\item कुछ हफ़्तों के अंदर खरीदने और बेचने का लक्ष्य रखें ताकि कम समय में ज़्यादा मुनाफ़ा कमाया जा सके।
\end{enumerate}

\bigskip


\subsection*{Question 24}
निम्नलिखित में से कौन सा Bitcoin प्राप्त करने के तीन प्राथमिक तरीकों में से एक नहीं है?

\begin{enumerate}[label=\Alph*., leftmargin=*]
\item जीडब्ल्यू-5 जीडब्ल्यू-4 विशेष हार्डवेयर और ऊर्जा उपयोग के माध्यम से।
\item प्लेटफॉर्म या ब्रोकर से Bitcoin खरीदना।
\item सरकारी अनुदान के रूप में Bitcoin प्राप्त करना, जो कारोबार के विकास के लिए है।
\item Exchange में Bitcoin प्राप्त करना (माल या सेवाओं के लिए)।
\end{enumerate}

\bigskip


\subsection*{Question 25}
Bitcoin खरीदते समय, इसकी कीमत में उतार-चढ़ाव को संभालने के लिए एक सुझाया गया तरीका क्या है?

\begin{enumerate}[label=\Alph*., leftmargin=*]
\item धीरे-धीरे नियमित अंतराल पर Bitcoin जमा करते रहें और लंबे समय तक का नजरिया बनाए रखें।
\item केवल उन दिनों Bitcoin खरीदें जब इसकी कीमत तेज़ी से बढ़ गई हो।
\item कंपनी का पूरा खजाना एक ही बार में झोंक दो।
\item छोटे समय के लिए मार्केट के चरम स्तर का अनुमान लगाने की कोशिश करो और तुरंत बेच दो।
\end{enumerate}

\bigskip


\subsection*{Question 26}
छोटी शुरुआत में Bitcoin की खरीदारी क्यों सुझाई जाती है व्यापारों के लिए?

\begin{enumerate}[label=\Alph*., leftmargin=*]
\item यह कंपनी को सभी क्रिप्टोकरेंसी स्कैम्स से बचाता है।
\item यह कानूनी तौर पर कंपनी को टैक्स के दायित्वों से बचाता है।
\item यह तुरंत व्यापार को अधिकतम मुनाफ़ा सुनिश्चित करता है।
\item यह उन्हें बड़ी रकम लगाने से पहले हाथों-हाथ अनुभव और समझ हासिल करने का मौका देता है।
\end{enumerate}

\bigskip


\subsection*{Question 27}
व्यापार के लिए Bitcoin खरीदने के लिए प्लेटफॉर्म या ब्रोकर चुनते समय एक महत्वपूर्ण बात क्या है?

\begin{enumerate}[label=\Alph*., leftmargin=*]
\item एक ऐसा प्रोवाइडर ढूंढना जो किसी भी कस्टमर सपोर्ट को देने से साफ़ मना कर दे।
\item प्रदाता की प्रतिष्ठा, सुरक्षा उपाय, और लेखांकन के लिए लेन-देन इतिहास निर्यात करने की क्षमता।
\item प्लेटफॉर्म को सिर्फ़ ग़ैर-मशहूर क्रिप्टोकरेंसी ही सपोर्ट करने की पक्की व्यवस्था करना।
\item एक ऐसी सेवा का इस्तेमाल करना जो Bitcoin को सेल्फ-कस्टडी Wallet में निकालने पर रोक लगाती है।
\end{enumerate}

\bigskip


\subsection*{Question 28}
Bitcoin प्राप्त करते समय कस्टडी के बारे में एक सुझावित प्रथा क्या है?

\begin{enumerate}[label=\Alph*., leftmargin=*]
\item सभी Bitcoin को एक सार्वजनिक वेबसाइट पर स्टोर करें, एक साधारण पासवर्ड के साथ
\item हमेशा सुरक्षा उपायों की जाँच किए बिना एक ही कस्टोडियन पर भरोसा करो।
\item शुरुआती खरीदारी के बाद सेल्फ-कस्टडी सॉल्यूशन्स पर विचार करें, ताकि बिजनेस वास्तव में अपनी संपत्ति पर कंट्रोल रख सके।
\item किसी भी हालत में Exchange से Bitcoin को वापस नहीं निकालना।
\end{enumerate}

\bigskip


\subsection*{Question 29}
कोई व्यवसाय ग्राहकों से Bitcoin भुगतान स्वीकार करना क्यों चुन सकता है?

\begin{enumerate}[label=\Alph*., leftmargin=*]
\item चेकआउट प्रक्रिया को धीमा और अधिक जटिल बनाना।
\item बढ़ते हुए उपयोगकर्ता आधार का लाभ उठाने, संभावित रूप से कम लेनदेन शुल्क, तथा अपने खजाने में एक मूल्यवान परिसंपत्ति रखने के लिए।
\item यह गारंटी देना कि सभी ग्राहक सूचीबद्ध मूल्य से अधिक भुगतान करेंगे।
\item यह सुनिश्चित करना कि कभी भी कोई कर या विनियामक आवश्यकताएं लागू न हों।
\end{enumerate}

\bigskip


\subsection*{Question 30}
जी, पॉइंट ऑफ सेल (POS) पर Lightning Network पेमेंट्स का मर्चेंट्स के लिए एक फायदा यह है कि इसमें **ज़ीरो ट्रांजैक्शन फीस** होती है। यानी, मर्चेंट्स को किसी भी अतिरिक्त चार्ज या कमीशन का भुगतान नहीं करना पड़ता, जिससे उनकी लागत कम होती है और मुनाफ़ा बढ़ता है। साथ ही, यह सिस्टम तेज़ और सुरक्षित होता है, जिससे ग्राहकों को भी आसानी होती है।

\begin{enumerate}[label=\Alph*., leftmargin=*]
\item बिजली की तरह तेज़ भुगतान, तुरंत सेटल हो जाते हैं और बिचौलियों की ज़रूरत नहीं पड़ती, जिससे फीस और देरी दोनों कम हो जाती है।
\item बिजली के भुगतान को पुष्टि करने में हफ्तों लगते हैं और भारी कागजी कार्रवाई की जरूरत होती है।
\item बिजली भुगतान तभी होते हैं जब किसी केंद्रीय बैंक के अधिकारी द्वारा मंजूरी मिल जाए।
\item खरीदार बिक्री के बाद मनमर्जी से लाइटनिंग भुगतान को वापस ले सकता है।
\end{enumerate}

\bigskip


\subsection*{Question 31}
Bitcoin पेमेंट्स स्वीकार करने से किसी बिज़नेस की विज़िबिलिटी कैसे बढ़ सकती है?

\begin{enumerate}[label=\Alph*., leftmargin=*]
\item सरकार द्वारा बिज़नेस को मुफ्त में प्रचार करने की गारंटी देकर।
\item सभी दूसरे भुगतान के तरीकों पर रोक लगाकर और सभी ग्राहकों को Bitcoin का इस्तेमाल करने के लिए मजबूर करके।
\item ग्राहकों को भुगतान करने से पहले कई ऑनलाइन खातों के लिए रजिस्टर करने के लिए मजबूर करके।
\item बिज़नेस को BTCmap.org जैसे प्लेटफॉर्म पर लिस्ट करके और Bitcoin कम्युनिटी के मुंहज़ोर प्रचार से।
\end{enumerate}

\bigskip


\subsection*{Question 32}
एक बिज़नेस को अपनी ज़रूरतों के हिसाब से Bitcoin पेमेंट सॉल्यूशन कैसे चुनना चाहिए?

\begin{enumerate}[label=\Alph*., leftmargin=*]
\item बिना बिज़नेस की खास ज़रूरतों को ध्यान में रखे, सिर्फ अंदाज़े पर भरोसा करके।
\item सिर्फ सबसे बड़ा और सबसे महंगा समाधान इस्तेमाल करके, चाहे बिज़नेस का आकार कुछ भी हो।
\item लेन-देन की आवृत्ति, औसत भुगतान राशि, और यह ऑनलाइन है या व्यक्तिगत रूप से - जैसे कारकों को ध्यान में रखकर, और फिर उस श्रेणी के अनुरूप एक समाधान चुनकर।
\item हमेशा सबसे सरल समाधान चुनकर, चाहे लेन-देन की मात्रा या जटिलता कुछ भी हो।
\end{enumerate}

\bigskip


\subsection*{Question 33}
Bitcoin पेमेंट्स स्वीकारना शुरू करते समय छोटे स्तर पर शुरुआत करने का एक फायदा क्या है?

\begin{enumerate}[label=\Alph*., leftmargin=*]
\item यह बिज़नेस को तुरंत सिर्फ़ सेल्फ-कस्टडी सॉल्यूशंस इस्तेमाल करने पर मजबूर कर देता है।
\item यह बिज़नेस को बाद में प्रोवाइडर्स को अपग्रेड या बदलने से रोकता है।
\item इसके लिए कंपनी को प्रक्रिया समझने से पहले हर ग्राहक से Bitcoin स्वीकार करना जरूरी होता है।
\item यह व्यवसाय को कुछ शुरुआती लेन-देन से सीखने और धीरे-धीरे अधिक उन्नत समाधानों की ओर बढ़ने की अनुमति देता है, अगर ग्राहकों की मांग बढ़ती है।
\end{enumerate}

\bigskip


\subsection*{Question 34}
Lightning Network ने अपनी शुरुआत से लेकर अब तक यूजर एक्सपीरियंस को बेहतर बनाने के लिए कैसे विकास किया है?

\begin{enumerate}[label=\Alph*., leftmargin=*]
\item इसने स्प्लिसिंग, स्टैटिक Bolt12 इनवॉइस, और बेहतर लिक्विडिटी मैनेजमेंट जैसी फीचर्स पेश की हैं, जिससे ट्रांजैक्शन्स और स्मूद और रिलायबल हो गए हैं।
\item यह नए फीचर्स पर काम करना बंद कर चुका है, सिर्फ़ पुराने डिज़ाइन पर ही भरोसा कर रहा है।
\item अब हर भुगतान के लिए सरकारी मंजूरी की जरूरत होती है।
\item इसने लेन-देन को तेज़ करने के लिए सभी सुरक्षा फ़ीचर्स हटा दिए हैं।
\end{enumerate}

\bigskip


\subsection*{Question 35}
Lightning Network पर नोड्स चलाने वाले प्रतिभागियों में क्या ट्रेंड देखने को मिल रहा है?

\begin{enumerate}[label=\Alph*., leftmargin=*]
\item व्यक्तिगत शौकीनों ने पूरी तरह से सभी पेशेवर ऑपरेटरों को बदल दिया है।
\item बड़े बैंक अब सभी लाइटनिंग नोड्स को केंद्रीय रूप से संचालित करते हैं।
\item पेशेवर नोड ऑपरेटर तेजी से लिक्विडिटी प्रदान कर रहे हैं, जबकि व्यक्तिगत यूजर्स "एज" नोड्स पर ध्यान केंद्रित कर रहे हैं, जिससे नेटवर्क की संरचना बदल रही है।
\item नेटवर्क ने बड़ी कंपनियों को ही अनुमति देने के लिए व्यक्तिगत नोड्स को बंद कर दिया है।
\end{enumerate}

\bigskip


\subsection*{Question 36}
Lightning Network की बेहतर कार्यक्षमता और स्पीड पारंपरिक भुगतान तरीकों की तुलना में कैसी है?

\begin{enumerate}[label=\Alph*., leftmargin=*]
\item इसे कई मंजूरी देने वाली संस्थाओं की जरूरत होती है, बिल्कुल वैसे ही जैसे बैंक ट्रांसफर में होता है।
\item यह पारंपरिक बैंकिंग समय से मेल खाता है, और अक्सर सेटल होने में 2-5 दिन लगते हैं।
\item यह काफी तेज़ है, अक्सर लेन-देन एक सेकंड से भी कम समय में पूरा कर देता है, जिसमें कम फीस लगती है और कई बिचौलियों पर निर्भरता नहीं होती।
\item यह पारंपरिक कार्ड भुगतान से धीमा है और इसकी फीस भी ज्यादा है।
\end{enumerate}

\bigskip


\subsection*{Question 37}
Lightning Network कैसे ऑनलाइन कॉमर्स और डिजिटल सेवाओं के भविष्य को आकार देने में मदद कर रहा है?

\begin{enumerate}[label=\Alph*., leftmargin=*]
\item बिना रुकावट, तुरंत और बिना अनुमति के माइक्रोट्रांजैक्शन को सक्षम करके, नए बिजनेस मॉडल और कमाई के नए रास्ते खोल दिए हैं।
\item डिजिटल सामान को बेचना मुश्किल बना देने से, ज़्यादा फीस की वजह से।
\item ऑनलाइन कोई नया बिजनेस मॉडल उभरने से रोककर।
\item व्यापारियों को पहले जैसी ही धीमी, महंगी भुगतान प्रणाली का उपयोग करने के लिए मजबूर करके।
\end{enumerate}

\bigskip


\subsection*{Question 38}
Bitcoin पेमेंट्स इकोसिस्टम का एक बड़ा फायदा यह है कि आप आसानी से अलग-अलग सॉल्यूशंस के बीच स्विच कर सकते हैं।

\begin{enumerate}[label=\Alph*., leftmargin=*]
\item प्रदाता बदलने के लिए सरकारी मंजूरी और लंबी कानूनी प्रक्रिया की जरूरत होती है।
\item जीडब्ल्यू-2 सॉल्यूशन्स प्रोप्राइटरी हैं और सिर्फ एक वेंडर तक लॉक्ड हैं।
\item यह प्रदाताओं को बदलना या टूल्स को अपग्रेड करना अपेक्षाकृत आसान है, क्योंकि Bitcoin के खुले, इंटरऑपरेबल मानकों की वजह से फंड ट्रांसफर बिना किसी रुकावट के हो जाता है।
\item एक बार प्रोवाइडर चुन लिया जाता है, तो बिज़नेस कभी दूसरे सॉल्यूशन पर नहीं जा सकता।
\end{enumerate}

\bigskip


\subsection*{Question 39}
कंपनी के अकाउंटिंग में Bitcoin के साथ काम करते समय सावधानीपूर्वक रिकॉर्ड रखना क्यों जरूरी है?

\begin{enumerate}[label=\Alph*., leftmargin=*]
\item क्योंकि Bitcoin लेनदेन को एक बार रिकॉर्ड करने के बाद उसका पता लगाना नामुमकिन हो जाता है।
\item क्योंकि Blockchain अपने आप कंपनी की तरफ से टैक्स फाइल कर देता है।
\item क्योंकि Bitcoin लेनदेन पर कोई भी लेखा मानक लागू नहीं होता है।
\item क्योंकि हर Bitcoin लेन-देन से टैक्स योग्य घटना हो सकती है, और मुनाफ़ा, नुकसान और टैक्स दायित्वों की गणना के लिए सटीक रिकॉर्ड्स की ज़रूरत होती है।
\end{enumerate}

\bigskip


\subsection*{Question 40}
Bitcoin को डिजिटल परिसंपत्ति के रूप में मानने और लेखांकन में इसे मुद्रा के रूप में मानने के बीच एक मुख्य अंतर क्या है?

\begin{enumerate}[label=\Alph*., leftmargin=*]
\item जब इसे डिजिटल परिसंपत्ति के रूप में माना जाता है, तो यह सभी करों से पूरी तरह मुक्त होती है।
\item एक डिजिटल परिसंपत्ति के रूप में, प्रत्येक लेनदेन के लागत आधार और लाभ/हानि को ट्रैक किया जाना चाहिए, जबकि यदि इसे मुद्रा की तरह माना जाता, तो पुनर्मूल्यांकन सरल और कम बार होता।
\item जब इसे मुद्रा की तरह व्यवहार किया जाता है, तो लेन-देन के रिकॉर्ड की आवश्यकता नहीं होती।
\item इसे मुद्रा की तरह व्यवहार करने के लिए तत्काल एक नए लेखांकन सॉफ्टवेयर पर स्विच करने की आवश्यकता है।
\end{enumerate}

\bigskip


\subsection*{Question 41}
एक बिज़नेस Bitcoin को अपने बैलेंस शीट पर लॉन्ग-टर्म इन्वेस्टमेंट के तौर पर क्यों रखना चाहेगा?

\begin{enumerate}[label=\Alph*., leftmargin=*]
\item क्योंकि इससे करेंसी से जुड़े सभी टैक्स तुरंत खत्म हो जाते हैं।
\item क्योंकि यह अंतर्राष्ट्रीय लेखा मानकों द्वारा आवश्यक है।
\item क्योंकि यह मूल्य संचय का काम कर सकता है, समय के साथ इसकी कीमत बढ़ सकती है और यह महंगाई से बचाव का ज़रिया भी हो सकता है।
\item क्योंकि Bitcoin में खुद ही बिना किसी जोखिम के गारंटीड मुनाफा मिलता है।
\end{enumerate}

\bigskip


\subsection*{Question 42}
Bitcoin से सभी लेन-देन का CSV एक्सपोर्ट रखने की सलाह क्यों दी जाती है?

\begin{enumerate}[label=\Alph*., leftmargin=*]
\item क्योंकि CSV फाइलों से कॉस्ट बेसिस ट्रैक करने की जरूरत नहीं रह जाती।
\item यह हर लेन-देन का साफ़ और व्यवस्थित रिकॉर्ड देता है, जिससे अकाउंटेंट्स के लिए इन्हें फाइनेंशियल रिपोर्ट्स और टैक्स फाइलिंग में शामिल करना आसान हो जाता है।
\item क्योंकि डेटा एक्सपोर्ट करने से आपका टैक्स रिटर्न बिना किसी रिव्यू के अपने आप पूरा हो जाता है।
\item क्योंकि कागज़ की रसीदें ही कानूनी तौर पर मान्य लेखा दस्तावेज़ हैं।
\end{enumerate}

\bigskip


\subsection*{Question 43}
अधिकांश न्यायक्षेत्रों में, Bitcoin से निपटने के दौरान कर योग्य घटना क्या होती है?

\begin{enumerate}[label=\Alph*., leftmargin=*]
\item बिना किसी लेनदेन के Bitcoin को अनिश्चित काल तक धारण करना।
\item Bitcoin को बेचना, व्यापार करना या भुगतान के रूप में उपयोग करना अक्सर कर योग्य घटना को जन्म देता है।
\item Bitcoin की कीमत ऑनलाइन देखना।
\item अपने Bitcoin Wallet या कुंजियों का बैकअप लेना।
\end{enumerate}

\bigskip


\subsection*{Question 44}
Bitcoin आपको बिना किसी बिचौलिए के क्या करने की अनुमति देता है?

\begin{enumerate}[label=\Alph*., leftmargin=*]
\item व्यक्तिगत डेटा स्टोर करें
\item बैंक से लोन लें
\item इंटरनेट पर Exchange का मूल्य
\item केंद्रीकृत एप्लिकेशन बनाएं
\end{enumerate}

\bigskip


\subsection*{Question 45}
Lightning Network का उपयोग करने वाले व्यापारियों के लिए एक बड़ा फायदा क्या है?

\begin{enumerate}[label=\Alph*., leftmargin=*]
\item कम लेन-देन शुल्क
\item व्यापारिक करों का उन्मूलन
\item स्वतः बिक्री में वृद्धि
\item बैंकों द्वारा गारंटीकृत भुगतान
\end{enumerate}

\bigskip


\subsection*{Question 46}
ऑस्ट्रियन स्कूल "टाइम प्रिफरेंस" को कैसे परिभाषित करता है?

\begin{enumerate}[label=\Alph*., leftmargin=*]
\item Supply और मांग को संतुलित करने की क्षमता
\item अभी खपत करने की पसंद बजाय बाद में
\item छोटे समय के खर्चे को कम करने की चाहत
\item फौरन निवेश का बेहतरीन इंतज़ाम
\end{enumerate}

\bigskip


\subsection*{Question 47}
मुद्रा (करेंसी) का मुख्य काम क्या है?

\begin{enumerate}[label=\Alph*., leftmargin=*]
\item भौतिक रूप में आंतरिक मूल्य को संग्रहित करने के लिए
\item एक बिचौलिया बनकर काम करना जो मूल्य के Exchange को सरल बनाता है
\item एक ही अधिकारी के ज़रिए सभी आर्थिक लेन-देन को केंद्रीकृत करना
\item बार्टर को पूरी तरह से सीधे सामान-के-बदले-सामान Exchange से बदलने के लिए
\end{enumerate}

\bigskip


\subsection*{Question 48}
क्यों एक अच्छी करेंसी के लिए कमी होना ज़रूरी है?

\begin{enumerate}[label=\Alph*., leftmargin=*]
\item यह सुनिश्चित करता है कि मुद्रा हमेशा सोने जैसी किसी वस्तु द्वारा समर्थित होनी चाहिए
\item यह सुनिश्चित करता है कि मुद्रा हर व्यक्ति को आसानी से उपलब्ध हो।
\item यह सुनिश्चित करता है कि केवल भौतिक संपत्ति ही मुद्रा के रूप में इस्तेमाल की जा सकती है
\item यह वैल्यू को बचाता है
\end{enumerate}

\bigskip


\subsection*{Question 49}
पारंपरिक भुगतान प्रणालियों का मुख्य काम क्या है?

\begin{enumerate}[label=\Alph*., leftmargin=*]
\item फंड ट्रांसफर को सिर्फ आमने-सामने नकद लेन-देन तक सीमित करना
\item अंतरराष्ट्रीय व्यापार को आसान बनाने के लिए
\item एकल वैश्विक मुद्रा बनाना जिसे केंद्रीय प्राधिकरण द्वारा नियंत्रित किया जाए
\item दो पक्षों के बीच धन हस्तांतरण को सक्षम करने के लिए
\end{enumerate}

\bigskip


\subsection*{Question 50}
पारंपरिक भुगतान प्रणालियों में सिक्कों के मुकाबले कागज़ के नोट का क्या बड़ा फायदा है?

\begin{enumerate}[label=\Alph*., leftmargin=*]
\item कीमती धातुओं के जरिए बढ़ाया गया आंतरिक मूल्य
\item हर लेन-देन की स्वचालित जाँच
\item पोर्टेबिलिटी और कार्यक्षमता में वृद्धि
\item नकली उत्पादों के जोखिमों का पूरी तरह से खात्मा
\end{enumerate}

\bigskip


\subsection*{Question 51}
क्रेडिट कार्ड लेनदेन सामान्यतः साधारण नकद लेनदेन से किस प्रकार भिन्न होता है?

\begin{enumerate}[label=\Alph*., leftmargin=*]
\item इसके लिए केवल मुद्रा के भौतिक Exchange की आवश्यकता होती है
\item इसमें कई मध्यस्थ शामिल होते हैं
\item यह अनिवार्य करता है कि दोनों पक्ष समान वस्तुओं का उत्पादन करें
\item यह बिना किसी मध्यस्थ की भागीदारी के तुरन्त होता है
\end{enumerate}

\bigskip


\subsection*{Question 52}
ऐसा कौन सा व्यवसाय है जिसके लिए एक साधारण CSV फाइल अकाउंटिंग के लिए काफी हो सकती है?

\begin{enumerate}[label=\Alph*., leftmargin=*]
\item छोटे व्यवसाय
\item प्रमुख वित्तीय संस्थान
\item बड़ी बहुराष्ट्रीय कंपनियाँ
\item हाई-फ़्रीक्वेंसी ट्रेडिंग प्लेटफ़ॉर्म्स
\end{enumerate}

\bigskip


\subsection*{Question 53}
कोर्स में बताए गए Bitcoin ट्रांजैक्शन CSV फाइल में निम्नलिखित में से कौन सा डिटेल साफ़ तौर पर दर्ज नहीं होता है?

\begin{enumerate}[label=\Alph*., leftmargin=*]
\item लेन-देन की तारीख
\item जीडब्ल्यू-1 से जुड़े पते
\item समतुल्य फिएट मूल्य
\item इस्तेमाल किया गया लिपि प्रकार
\end{enumerate}

\bigskip


\subsection*{Question 54}
हर Bitcoin लेन-देन के लिए कौन-कौन से मुख्य भुगतान विवरण दर्ज किए जाने चाहिए?

\begin{enumerate}[label=\Alph*., leftmargin=*]
\item तारीख, Exchange दर, Wallet प्रकार, और भौगोलिक स्थान
\item जीडब्ल्यू-1, नेटवर्क फीस, और मर्चेंट का जीडब्ल्यू-2
\item जीडब्ल्यू-6 ब्लॉक नंबर, जीडब्ल्यू-5, और इस्तेमाल किया गया स्क्रिप्ट
\item तारीख, रकम, समतुल्य फिएट मूल्य, और संबंधित पतों
\end{enumerate}

\bigskip


\subsection*{Question 55}
निम्नलिखित में से कौन सा सॉफ्टवेयर Bitcoin अकाउंटिंग सॉल्यूशन्स में से नहीं है?

\begin{enumerate}[label=\Alph*., leftmargin=*]
\item कॉइनट्रैकर
\item कोइनली
\item ज़ेनलेजर
\item गौरैया
\end{enumerate}

\bigskip


\subsection*{Question 56}
विशेष Bitcoin एकाउंटिंग टूल्स का मुख्य लक्ष्य क्या है?

\begin{enumerate}[label=\Alph*., leftmargin=*]
\item Bitcoin प्राइवेट कीज़ की सुरक्षा बढ़ाने के लिए
\item Bitcoin पेमेंट स्वीकार करने के लिए PoS सिस्टम सेट अप करने के लिए
\item Bitcoin से जुड़े डेटा के स्वचालित संग्रह और प्रसंस्करण को सरल बनाने के लिए
\item व्यवसायों के लिए Bitcoin को सीधे डिसेंट्रलाइज्ड एक्सचेंजों पर ट्रेडिंग करने में आसानी लाने के लिए
\end{enumerate}

\bigskip


\subsection*{Question 57}
एक नॉन-कस्टोडियल Wallet स्टार्टर यूजर्स के लिए फंक्शनैलिटी में कैसे अलग होता है?

\begin{enumerate}[label=\Alph*., leftmargin=*]
\item यह बिज़नेस मालिक को प्राइवेट कीज़ का पूरा कंट्रोल देता है
\item यह सारे लेन-देन को एक तीसरे पक्ष की सेवा के जरिए संभालता है
\item इसे किसी बैकअप प्रक्रिया या सुरक्षा उपायों की जरूरत नहीं है।
\item यह अपने आप बिटकॉइन को फिएट करेंसी में बदल देता है
\end{enumerate}

\bigskip


\subsection*{Question 58}
यह क्यों जरूरी है कि हर Bitcoin ट्रांजैक्शन के वक्त फिएट वैल्यू को ट्रैक किया जाए?

\begin{enumerate}[label=\Alph*., leftmargin=*]
\item सटीक बहीखाता और टैक्स अनुपालन सुनिश्चित करने के लिए
\item नेटवर्क द्वारा लगाए गए लेन-देन शुल्क को कम करने के लिए
\item Wallet बैलेंस को अपने आप बढ़ाने के लिए
\item Wallet से सीधे रियल-टाइम मार्केट ट्रेडिंग को सक्षम करने के लिए
\end{enumerate}

\bigskip


\subsection*{Question 59}
एक फायदा यह है कि कस्टोडियल Wallet का इस्तेमाल करने से Starter प्रोफाइल वालों को अपने फंड्स का मैनेजमेंट करने की चिंता नहीं रहती, क्योंकि यह सुरक्षित तरीके से थर्ड पार्टी के ज़रिए होल्ड किया जाता है।

\begin{enumerate}[label=\Alph*., leftmargin=*]
\item प्राइवेट कीज़ पर बढ़ा हुआ कंट्रोल
\item लेन-देन पर पूर्ण नियंत्रण के साथ बेहतर गोपनीयता
\item टेक्निकल ज़िम्मेदारियाँ कम हो गईं
\item लेन-देन शुल्क को खत्म करना
\end{enumerate}

\bigskip


\subsection*{Question 60}
स्टार्टर प्रोफाइल में कौन से दो Wallet प्रकार माने जाते हैं?

\begin{enumerate}[label=\Alph*., leftmargin=*]
\item कस्टोडियल और नॉन-कस्टोडियल लाइटनिंग वॉलेट्स
\item जीडब्ल्यू-4 और पेपर जीडब्ल्यू-5 वॉलेट्स
\item हार्डवेयर और ब्रेन लाइटनिंग वॉलेट्स
\item जीडब्ल्यू-3 और जीडब्ल्यू-2 स्टोरेज जीडब्ल्यू-1 वॉलेट्स
\end{enumerate}

\bigskip


\subsection*{Question 61}
स्टार्टर प्रोफाइल के लिए Bitcoin पेमेंट स्वीकार करने की प्राथमिक हार्डवेयर आवश्यकता क्या है?

\begin{enumerate}[label=\Alph*., leftmargin=*]
\item एक जीडब्ल्यू-1 जीडब्ल्यू-2 जीडब्ल्यू-3
\item एक स्मार्टफोन
\item एक स्व-होस्टेड सर्वर
\item एक जीडब्ल्यू-4
\end{enumerate}

\bigskip


\subsection*{Question 62}
एसेंशियल प्रोफाइल पॉइंट ऑफ सेल (PoS) पर कर्मचारियों को भुगतान प्रोसेस करने में मदद के लिए क्या इस्तेमाल करता है?

\begin{enumerate}[label=\Alph*., leftmargin=*]
\item पीओएस ऐप पर आइटम की कीमतों के साथ एक मेनू बनाना
\item हर लेन-देन के लिए एक कागज़ी Invoice सिस्टम
\item पॉइंट ऑफ सेल (PoS) से लेन-देन के लिए एक स्वचालित लेखा रिकॉर्ड सिस्टम
\item एक मैनुअल जीडब्ल्यू-0 बुक जो सेल्स रिकॉर्ड करने के लिए है
\end{enumerate}

\bigskip


\subsection*{Question 63}
स्विस Bitcoin के बजाय कौन सा विकल्प Essential प्रोफाइल के लिए Bitcoin भुगतान स्वीकार करने में 100% कस्टोडियल है?

\begin{enumerate}[label=\Alph*., leftmargin=*]
\item बिटकॉइनाइज़ पीओएस
\item नोड खोलो
\item फीनिक्स जीडब्ल्यू-1
\item बीटीसीपे सर्वर
\end{enumerate}

\bigskip


\subsection*{Question 64}
कौन सा सीन Bitcoin पेमेंट्स के लिए एसेंशियल प्रोफाइल के इस्तेमाल को सबसे अच्छे तरीके से दिखाता है?

\begin{enumerate}[label=\Alph*., leftmargin=*]
\item एक स्वतंत्र फ्रीलांसर जो अपने इनवॉइस के भुगतान के लिए बिटकॉइन स्वीकार करता है
\item एक रेस्तरां जो महीने में कुछ Bitcoin भुगतान स्वीकार करता है
\item एक बहुराष्ट्रीय कंपनी जो रोज़ हज़ारों लेन-देन संभालती है
\item एक हाई-फ्रीक्वेंसी Bitcoin ट्रेडिंग डेस्क
\end{enumerate}

\bigskip


\subsection*{Question 65}
स्विस Bitcoin Pay सॉल्यूशन की कौन सी खासियत इसे 'एसेंशियल प्रोफाइल' वाले छोटे व्यवसायों के लिए खास तौर पर उपयुक्त बनाती है?

\begin{enumerate}[label=\Alph*., leftmargin=*]
\item एक मल्टी-सिग्नेचर जीडब्ल्यू-1 मैनेजमेंट
\item एडवांस्ड मल्टी-करेंसी ट्रेडिंग फीचर्स
\item एक यूजर-फ्रेंडली पॉइंट ऑफ सेल (PoS) ऐप
\item पूरा Blockchain एनालिटिक्स इंटीग्रेशन
\end{enumerate}

\bigskip


\subsection*{Question 66}
Essential प्रोफाइल बिज़नेस के लिए सबसे बढ़िया पेमेंट सॉल्यूशन कौन सा है?

\begin{enumerate}[label=\Alph*., leftmargin=*]
\item स्विस जीडब्ल्यू-0 पे
\item बिटकॉइनाइज़ पीओएस
\item ओपन नोड
\item बीटीसीपे सर्वर
\end{enumerate}

\bigskip


\subsection*{Question 67}
व्यवसायों के लिए कौन सा वैकल्पिक समाधान बताया गया है जो एक ऑल-इन-वन ऑनलाइन रिटेल वातावरण पसंद करते हैं?

\begin{enumerate}[label=\Alph*., leftmargin=*]
\item फीनिक्स
\item बी-बॉप
\item स्विस जीडब्ल्यू-0 पे
\item ओपन नोड
\end{enumerate}

\bigskip


\subsection*{Question 68}
टूल्स जैसे कि Zaprite और Musqet एक प्रोफेशनल पेमेंट सॉल्यूशन में क्या भूमिका निभाते हैं?

\begin{enumerate}[label=\Alph*., leftmargin=*]
\item वे बेसिक पॉइंट ऑफ सेल (POS) सिस्टम की तरह काम करते हैं
\item वे लेन-देन डेटा को एन्क्रिप्ट करने और प्राइवेट कीज़ को मैनेज करने पर ध्यान देते हैं
\item वे मल्टी-करेंसी कन्वर्ज़न और फिएट डिस्बर्समेंट संभालते हैं।
\item वे चेकआउट का अनुभव बेहतर बनाते हैं
\end{enumerate}

\bigskip


\subsection*{Question 69}
BTC Pay Server को Professional Environment में कैसे Deploy किया जा सकता है?

\begin{enumerate}[label=\Alph*., leftmargin=*]
\item समर्पित हार्डवेयर उपकरणों पर विशेष रूप से
\item ऑन-प्रिमाइसेस या क्लाउड होस्टिंग के जरिए
\item एक मोबाइल ऐप के रूप में
\item एक पूरी तरह से आउटसोर्स किए गए तीसरे पक्ष की सेवा के माध्यम से
\end{enumerate}

\bigskip


\subsection*{Question 70}
कौन सा ओपन-सोर्स सॉफ्टवेयर प्रोफेशनल Bitcoin पेमेंट सिस्टम के लिए अनुकूलित है?

\begin{enumerate}[label=\Alph*., leftmargin=*]
\item खुला नोड
\item बिटकॉइनाइज़ पीओएस (PoS)
\item बीटीसी पे सर्वर
\item स्विस जीडब्ल्यू-1 पे
\end{enumerate}

\bigskip


\subsection*{Question 71}
पेशेवर प्रोफाइल वाले बिज़नेस आमतौर पर कौन-सी एडवांस्ड फीचर्स की मांग करते हैं?

\begin{enumerate}[label=\Alph*., leftmargin=*]
\item सरल Wallet एकीकरण और QR कोड जनरेशन
\item बिना किसी रिपोर्टिंग विकल्प के स्वचालित रूप से फिएट में परिवर्तन
\item मूल लेन-देन रिकॉर्डिंग और मानक रसीदें
\item एडवांस्ड रिपोर्टिंग डैशबोर्ड्स और मल्टीपल एडमिनिस्ट्रेटिव रोल्स
\end{enumerate}

\bigskip


\subsection*{Question 72}
एंटरप्राइज़ Bitcoin पेमेंट्स में LSP क्या भूमिका निभाता है?

\begin{enumerate}[label=\Alph*., leftmargin=*]
\item लाइटनिंग Wallet के प्राइवेट कीज को स्टोर करें
\item लाइटनिंग पर मैन्युअली रिफंड प्रोसेस करें
\item प्रिंट लाइटनिंग इनवॉइस
\item लाइटनिंग पर भारी मात्रा में लेन-देन को प्रबंधित करें
\end{enumerate}

\bigskip


\subsection*{Question 73}
एंटरप्राइज़ सिस्टम कैसे अपने आप इन्वेंटरी अपडेट करते हैं?

\begin{enumerate}[label=\Alph*., leftmargin=*]
\item वर्कफ़्लो ट्रिगर्स के माध्यम से
\item नियोजित रिपोर्ट्स
\item दिन के अंत में अपलोड्स
\item मैन्युअल अपडेट्स
\end{enumerate}

\bigskip


\subsection*{Question 74}
एंटरप्राइज़ Bitcoin सिस्टम और ERP सॉफ़्टवेयर के बीच रियल-टाइम डेटा सिंक कैसे होता है?

\begin{enumerate}[label=\Alph*., leftmargin=*]
\item एपीआई (API)
\item CSV फाइलें
\item बैच अपलोड
\item ईमेल रिपोर्ट्स
\end{enumerate}

\bigskip


\subsection*{Question 75}
एंटरप्राइज़ लेवल पर बिटकॉइन के लिए आमतौर पर कौन-कौन से एडवांस्ड सिक्योरिटी फीचर्स इस्तेमाल किए जाते हैं?

\begin{enumerate}[label=\Alph*., leftmargin=*]
\item सिंगल-सिग वाले 24 शब्दों वाला seed फ्रेज़
\item मल्टी-सिग्नेचर टाइमलॉक के साथ
\item सिंगल-सिग (Single-sig) जो BIP39 passphrase प्रोटेक्शन के साथ है
\item मोबाइल लाइटनिंग Wallet के लिए पासवर्ड सुरक्षा
\end{enumerate}

\bigskip


\subsection*{Question 76}
एंटरप्राइज़-लेवल Bitcoin पेमेंट इम्प्लीमेंटेशन का मुख्य उद्देश्य क्या है?

\begin{enumerate}[label=\Alph*., leftmargin=*]
\item संगठन के मुख्य प्रक्रियाओं में Bitcoin भुगतानों को पूरी तरह से शामिल करने के लिए
\item मार्केटिंग के लिए Bitcoin का इस्तेमाल करने के लिए
\item Bitcoin को छोटे पैमाने पर आज़माने के लिए
\item Bitcoin को अन्य सिस्टम्स में इंटीग्रेट किए बिना स्वीकार करना
\end{enumerate}

\bigskip


\newpage
\section*{Answer Key}

\begin{enumerate}
\item \textbf{Question 1:} C
\\ \textit{कोई भी चीज़ मुफ्त नहीं होती। पारंपरिक मुद्रा प्रणालियों में बने भुगतान नेटवर्क को कार्ड कंपनियों और अधिग्रहण बैंकों जैसे विभिन्न विक्रेताओं द्वारा संचालित किया जाता है, जो हर लेन-देन के लिए एक निश्चित शुल्क और/या प्रतिशत शुल्क वसूलते हैं।}

\item \textbf{Question 2:} D
\\ \textit{एक ऐसी मुद्रा जिसकी उपलब्धता सीमित हो, वह समय के साथ अपनी क्रय शक्ति बनाए रखती है। अगर यह दुर्लभता न हो, तो ज़रूरत से ज़्यादा मुद्रा छापने से मुद्रास्फीति हो सकती है, जिससे मुद्रा का मूल्य कम हो जाता है।}

\item \textbf{Question 3:} C
\\ \textit{कागज़ के नोट कीमती धातुओं से बने सिक्कों के मुकाबले हल्के और बनाने व ढोने में आसान थे। यह पोर्टेबिलिटी और मानकीकरण ने लेन-देन को सरल बना दिया और अर्थव्यवस्थाओं को स्थानीय क्षेत्रों से आगे बढ़ने में मदद की।}

\item \textbf{Question 4:} C
\\ \textit{केंद्रीकृत सिस्टम मनमर्जी से पैसा छाप सकते हैं, अक्सर बिना किसी पारदर्शी निगरानी के। इससे मुद्रास्फीति, जटिल इंटरबैंक सेटलमेंट्स और संभावित गलत प्रबंधन के कारण भरोसा कम हो सकता है।}

\item \textbf{Question 5:} C
\\ \textit{कार्ड पेमेंट्स में अक्सर कई बिचौलिये शामिल होते हैं (एक्वायरर्स, इश्यूअर्स, कार्ड नेटवर्क्स), इन्हें अथॉराइज़ेशन और सेटलमेंट के चरणों की ज़रूरत होती है, और मर्चेंट को नेट फंड्स मिलने में 48 घंटे से लेकर कई दिन तक का समय लग सकता है।}

\item \textbf{Question 6:} D
\\ \textit{Bitcoin एक विकेंद्रीकृत नेटवर्क की तरह काम करता है, जहाँ बिना किसी अनुमति के कोई भी हिस्सा ले सकता है। लेन-देन को नोड्स और माइनर्स द्वारा वैलिडेट किया जाता है, न कि किसी केंद्रीय प्राधिकरण द्वारा।}

\item \textbf{Question 7:} C
\\ \textit{21 मिलियन की अधिकतम Supply सीमा होने से, Bitcoin की कमी Hard कोडेड है, जो इसे मुद्रास्फीति वाली करेंसी से बिल्कुल अलग बनाता है जहां Supply को मनमर्जी से बढ़ाया जा सकता है।}

\item \textbf{Question 8:} C
\\ \textit{Bitcoin का Blockchain एक हैक-प्रूफ Ledger है। एक बार जब कोई लेन-देन कन्फर्म हो जाता है और ब्लॉक में शामिल हो जाता है, तो वह परमानेंट रिकॉर्ड का हिस्सा बन जाता है, जिससे फाइनैलिटी सुनिश्चित होती है।}

\item \textbf{Question 9:} D
\\ \textit{Bitcoin का ग्लोबल, इंटरनेट-नेटिव डिज़ाइन इसे सिर्फ एक वैल्यू स्टोर के रूप में नहीं, बल्कि Exchange के माध्यम के तौर पर भी काम करने देता है। यह व्यवसायों को दुनिया भर में तेज़, बिना अनुमति के और अंतिम लेन-देन सेटलमेंट की सुविधा देता है।}

\item \textbf{Question 10:} D
\\ \textit{Lightning Network एक "दूसरी पीढ़ी का Layer" प्रोटोकॉल है जो प्रतिभागियों को भुगतान चैनल बनाने, off-chain लेनदेन करने और अंतिम शुद्ध शेष को सिर्फ Bitcoin Blockchain पर सेटल करने की अनुमति देता है। यह डिज़ाइन तेज़ गति और कम लागत हासिल करता है।}

\item \textbf{Question 11:} B
\\ \textit{जब लाइटनिंग ट्रांजैक्शन off-chain पर होते हैं, तो सिर्फ फाइनल चैनल स्टेट्स Bitcoin Blockchain पर सेटल होती हैं। इससे कंजेशन कम होता है और फीस भी घट जाती है क्योंकि बहुत सारे On-Chain ट्रांजैक्शन करने से बचा जाता है।}

\item \textbf{Question 12:} B
\\ \textit{लाइटनिंग के साथ, एक बार पेमेंट नेटवर्क के ज़रिए रूट और सेटल हो जाता है, तो वो फाइनल होता है। यह पारंपरिक पेमेंट सिस्टम में देखे जाने वाले चार्जबैक के मामलों को खत्म कर देता है।}

\item \textbf{Question 13:} B
\\ \textit{जैसे ईमेल ने बिचौलियों को हटाकर और दुनिया भर में किसी को भी संदेश भेजने की सुविधा देकर संचार में क्रांति ला दी, वैसे ही Lightning Network लेन-देन को सरल और तेज़ बनाता है, जिससे मूल्य हस्तांतरण ईमेल की तरह ही आसान और सुलभ हो जाता है।}

\item \textbf{Question 14:} A
\\ \textit{ऑस्ट्रियन इकोनॉमिक्स में, मुनाफ़ा यह दिखाता है कि कंपनी ऐसे सामान या सेवाएँ बना रही है जिन्हें ग्राहक उत्पादन लागत से ज़्यादा कीमत पर वैल्यू देते हैं। यह नतीजा बताता है कि संसाधनों का कुशल इस्तेमाल हो रहा है और यह बाज़ार की ज़रूरतों के साथ मेल खा रहा है।}

\item \textbf{Question 15:} C
\\ \textit{भविष्य की स्थितियाँ अंदाज़ा लगाना मुश्किल होता है। कुछ पूँजी रिज़र्व में रखकर, कोई भी व्यवसाय नए मौकों का फायदा उठा सकता है, अचानक आने वाली मुश्किलों से निपट सकता है, और बाहरी वित्त पर निर्भर हुए बिना लंबे समय के प्रोजेक्ट्स में निवेश कर सकता है।}

\item \textbf{Question 16:} B
\\ \textit{वो मछुआरा जो कुछ दिनों की मछलियाँ बचाकर एक भाला बनाता है, वह अपनी तात्कालिक खपत को भविष्य के फायदे के लिए टाल रहा है। यह उदाहरण दिखाता है कि बचत (पूँजी संचय) कैसे भविष्य में अधिक दक्षता और ज्यादा उत्पादन की ओर ले जाती है।}

\item \textbf{Question 17:} C
\\ \textit{ऑस्ट्रियाई राजधानी असली बचत और सोच-समझकर निवेश से पैदा होती है। कर्ज़े पर बनी पूँजी, जो आसानी से लोन के ज़रिए बन जाती है, व्यापारियों को गुमराह कर सकती है और बेकार निवेश की वजह बन सकती है, क्योंकि यह असली बचत या ग्राहकों की पसंद को नहीं दिखाती।}

\item \textbf{Question 18:} C
\\ \textit{Bitcoin की कमी, वैश्विक तरलता, और पारंपरिक बैंकिंग प्रणालियों से स्वतंत्रता इसे एक ट्रेजरी संपत्ति के रूप में आकर्षक बना सकती है। हालांकि इसमें जोखिम और अस्थिरता शामिल है, फिर भी व्यवसाय इसे मुद्रास्फीति के खिलाफ एक बचाव और दीर्घकालिक मूल्य संचय के रूप में देख सकते हैं।}

\item \textbf{Question 19:} A
\\ \textit{जब बड़े संस्थान Bitcoin को एक वैध निवेश के रूप में मान्यता देते हैं, तो यह बाजार को संकेत देता है कि यह अटकलबाजी से आगे बढ़ चुका है। स्पॉट Bitcoin ETFs जैसे उत्पाद व्यवसायों के लिए एक्सपोजर हासिल करना आसान बनाते हैं, जिससे मुख्यधारा में स्वीकृति और बढ़ती है।}

\item \textbf{Question 20:} A
\\ \textit{Bitcoin एक विकासशील संपत्ति है, जिसकी कीमत अटकलों, मनोदशा और अलग-अलग अपनाने की दरों से प्रभावित होती है। समय के साथ, जब अधिक स्थिर, दीर्घकालिक धारक भाग लेंगे, तो अस्थिरता कम हो सकती है, लेकिन यह शुरुआती चरण के विकास की एक खासियत बनी रहेगी।}

\item \textbf{Question 21:} A
\\ \textit{छुपी हुई मुद्रास्फीति तब होती है जब पैसा Supply बढ़ता है, जिससे बचत की खरीदारी शक्ति कम हो जाती है। अगर कोई व्यवसाय सक्रिय रूप से निवेश नहीं करता या अपनी दौलत की सुरक्षा नहीं करता, तो उसे मुद्रा के इस चुपके से कमजोर होने के कारण मूल्य खोने का खतरा होता है।}

\item \textbf{Question 22:} A
\\ \textit{रियल एस्टेट में अक्सर लिक्विडिटी की कमी और लगातार खर्चे होते हैं, और यह कर्ज़-आधारित मौद्रिक प्रणाली में छुपी हुई मुद्रास्फीति के साथ कदम नहीं मिला पाता। Bitcoin की कमी और वैश्विक पहुंच एक अलग स्टोर-ऑफ-वैल्यू प्रस्ताव देती है जो लंबे समय में पारंपरिक संपत्तियों से बेहतर प्रदर्शन कर सकती है।}

\item \textbf{Question 23:} C
\\ \textit{क्योंकि Bitcoin की कीमत कम समय में उतार-चढ़ाव भरी हो सकती है, इसलिए बेहतर है कि आप ऐसे पैसे इस्तेमाल करें जिनकी आपको अभी जरूरत नहीं है और कम से कम चार साल तक होल्ड करें। इस तरह से उतार-चढ़ाव कम होता है और लंबे समय में कीमत बढ़ने की संभावना बेहतर होती है।}

\item \textbf{Question 24:} C
\\ \textit{अध्याय में Bitcoin प्राप्त करने के तीन मुख्य तरीकों का जिक्र है; इसे भुगतान के रूप में स्वीकार करना, खनन करना, या खरीदना। सरकारी अनुदान के रूप में Bitcoin प्राप्त करना, अधिग्रहण का मान्यता प्राप्त मानक तरीका नहीं है।}

\item \textbf{Question 25:} A
\\ \textit{समय के साथ नियमित खरीदारी (डॉलर-कॉस्ट एवरेजिंग) उतार-चढ़ाव को कम करने में मदद करती है। लंबी अवधि (जैसे, 4 साल) का नजरिया अल्पकालिक कीमतों में उछाल के प्रभाव को कम करता है।}

\item \textbf{Question 26:} D
\\ \textit{एक छोटा सा शुरुआती निवेश व्यवसायों को Bitcoin खरीदने, स्टोर करने और प्रबंधित करने के व्यावहारिक पहलुओं को सीखने में मदद करता है, बिना कोई बड़ा जोखिम उठाए।}

\item \textbf{Question 27:} B
\\ \textit{एक भरोसेमंद प्लेटफॉर्म की मजबूत प्रतिष्ठा होनी चाहिए, अच्छी सुरक्षा होनी चाहिए, और रिकॉर्ड-रखने के लिए टूल्स भी होने चाहिए। Bitcoin निकालने की क्षमता और अकाउंटिंग सॉफ्टवेयर के साथ इंटीग्रेशन भी बिजनेस ऑपरेशन्स के लिए बहुत जरूरी है।}

\item \textbf{Question 28:} C
\\ \textit{Bitcoin खरीदने के बाद, इसे सेल्फ-कस्टडी Wallet में ट्रांसफर करने से बिज़नेस को अपने फंड्स पर सीधा कंट्रोल मिलता है। हालांकि थर्ड-पार्टी कस्टडियन सुविधाजनक हो सकते हैं, लेकिन सेल्फ-कस्टडी Bitcoin के फाइनेंशियल सॉवरेन्टी के सिद्धांत के साथ मेल खाती है।}

\item \textbf{Question 29:} B
\\ \textit{Bitcoin को स्वीकार करने से खरीदारों का एक नया जनसांख्यिकीय समूह आकर्षित हो सकता है, भुगतान प्रसंस्करण शुल्क कम हो सकता है, और मूल्य का भंडार बन सकता है। यह वैश्विक पहुंच को भी बढ़ा सकता है, क्योंकि Bitcoin सीमाहीन है और चौबीसों घंटे काम करता है।}

\item \textbf{Question 30:} A
\\ \textit{Lightning Network, off-chain सेटलमेंट की मदद से लेन-देन को तेज़ और सस्ता बना देता है। दुकानदारों को पैसे लगभग तुरंत मिल जाते हैं, बहुत कम फीस के साथ, और बिना पारंपरिक भुगतान बिचौलियों पर निर्भर रहे।}

\item \textbf{Question 31:} D
\\ \textit{Bitcoin को स्वीकार करना अभी भी काफी कम है, जिससे Bitcoin कम्युनिटी का ध्यान आकर्षित हो सकता है। अपनी दुकान को विशेष डायरेक्टरीज़ में लिस्ट करना और कम्युनिटी-चालित प्रमोशन का फायदा उठाने से आप अलग दिख सकते हैं और नए ग्राहकों को आकर्षित कर सकते हैं।}

\item \textbf{Question 32:} C
\\ \textit{अलग-अलग बिज़नेस की अलग-अलग ज़रूरतें होती हैं। लेन-देन की मात्रा, भुगतान के तरीके (ऑनलाइन बनाम फिजिकल), और स्केलेबिलिटी जैसे फैक्टर्स को देखकर कंपनी सही Bitcoin पेमेंट सॉल्यूशन चुन सकती है।}

\item \textbf{Question 33:} D
\\ \textit{एक साधारण, शुरुआती स्तर के तरीके से शुरुआत करने से व्यवसाय को बिना बड़े जोखिम के व्यावहारिक अनुभव हासिल करने में मदद मिलती है। अगर Bitcoin भुगतानों में रुचि फायदेमंद साबित होती है, तो कंपनी बाद में अपने बुनियादी ढांचे को अपग्रेड कर सकती है और और अधिक उन्नत टूल्स को शामिल कर सकती है।}

\item \textbf{Question 34:} A
\\ \textit{समय के साथ, Lightning Network में तकनीकी सुधार हुए हैं जिससे यूजर्स के लिए घर्षण कम हुआ है। नई सुविधाएँ और बेहतर नोड प्रबंधन टूल्स ने एक सहज और यूजर-फ्रेंडली अनुभव में योगदान दिया है।}

\item \textbf{Question 35:} C
\\ \textit{पेशेवर संस्थाएँ ज़्यादा लिक्विडिटी और स्थिरता दे रही हैं, जबकि आम यूज़र्स सरल या होस्ट किए गए समाधानों पर निर्भर हैं, जिससे नेटवर्क और मज़बूत और सुलभ हो रहा है।}

\item \textbf{Question 36:} C
\\ \textit{off-chain लेनदेन को सक्षम करके और भुगतान को लगभग तुरंत अंतिम रूप देकर, Lightning पारंपरिक तरीकों पर गति और दक्षता का लाभ प्रदान करता है जो कई मध्यस्थों पर निर्भर करते हैं।}

\item \textbf{Question 37:} A
\\ \textit{बिजली की तरह तेज़ और कम खर्च में छोटी-छोटी रकम भेजने की क्षमता ने नए आइडियाज़ को बढ़ावा दिया है, जैसे कि "जितना इस्तेमाल करो, उतना भुगतान करो" कंटेंट, लगातार भुगतान की सुविधा, और मशीनों के बीच लेन-देन। यह तकनीक ऑनलाइन सेवाओं से पैसा कमाने के तरीकों को पूरी तरह बदल सकती है!}

\item \textbf{Question 38:} C
\\ \textit{चूंकि Bitcoin ओपन प्रोटोकॉल्स पर बना है, इसलिए व्यापारी अपने फंड्स को कहीं भी मूव कर सकते हैं या जैसे-जैसे उनका बिज़नेस बढ़ता है या ज़रूरतें बदलती हैं, नए सॉल्यूशंस अपना सकते हैं। यह लचीलापन यह सुनिश्चित करता है कि बिज़नेसेस बिना किसी वेंडर लॉक-इन के सबसे अच्छे टूल्स चुन सकें।}

\item \textbf{Question 39:} D
\\ \textit{Bitcoin में की गई हर खरीद, बिक्री या भुगतान पर टैक्स का असर हो सकता है। खर्च का आधार (cost basis), बिक्री मूल्य और समय की विस्तृत रिकॉर्डिंग रखने से सही और कानून के अनुसार टैक्स रिपोर्टिंग सुनिश्चित होती है।}

\item \textbf{Question 40:} B
\\ \textit{डिजिटल एसेट ट्रीटमेंट में आमतौर पर प्रत्येक लेनदेन पर पूंजीगत लाभ या हानि की गणना करना शामिल होता है, जिससे जटिलता बढ़ जाती है। अगर Bitcoin को मुद्रा की तरह हिसाब में लिया जाता, तो यह आसान होता; बस समय-समय पर फ़िएट मूल्य को अपडेट करें, जैसा कि अन्य मुद्राओं के लिए किया जाता है।}

\item \textbf{Question 41:} C
\\ \textit{कुछ कंपनियां Bitcoin को एक स्ट्रैटेजिक लॉन्ग-टर्म एसेट के रूप में देखती हैं जिसकी वैल्यू प्रिज़र्व या ग्रो हो सकती है, जो ट्रेडिशनल करेंसीज़ के डिवैल्यूएशन से बचाती है और उनके ट्रेजरी को डायवर्सिफाई करती है।}

\item \textbf{Question 42:} B
\\ \textit{एक कॉमन फॉर्मेट में डिजिटल रिकॉर्ड रखने से अकाउंटिंग सिस्टम में आसानी से जुड़ जाता है, ऑडिट ट्रेल को मेंटेन करने में मदद मिलती है, और टैक्स व कंप्लायंस रिपोर्टिंग एकदम सही होती है।}

\item \textbf{Question 43:} B
\\ \textit{आम तौर पर, कर योग्य घटना तब होती है जब कोई व्यवसाय अपने Bitcoin का निपटान करता है या उसका उपयोग करता है, न कि केवल उसे धारण करके। माल, सेवाओं या फिएट मुद्रा के लिए Bitcoin की प्रत्येक बिक्री या Exchange generate कर योग्य लाभ या हानि हो सकती है।}

\item \textbf{Question 44:} C
\\ \textit{Bitcoin, अपने पीयर-टू-पीयर प्रोटोकॉल के जरिए, बिना किसी बिचौलिए के इंटरनेट पर Exchange के मूल्य को सक्षम बनाता है।}

\item \textbf{Question 45:} A
\\ \textit{Lightning Network पर लेन-देन शुल्क बहुत कम है, जिससे यह व्यापारियों के लिए एक किफायती समाधान बन जाता है।}

\item \textbf{Question 46:} B
\\ \textit{समय प्राथमिकता (Time preference) का मतलब है तुरंत खर्च करने को चुनना बजाय भविष्य के फायदे के, जो बचत और निवेश के फैसलों को प्रभावित करता है।}

\item \textbf{Question 47:} B
\\ \textit{मुद्रा व्यापार को आसान बनाती है क्योंकि यह एक सामान्य माध्यम के रूप में काम करती है जो बार्टर प्रणाली में जटिल Exchange दरों के जाल को कम कर देती है, जिससे लेन-देन आसान हो जाता है।}

\item \textbf{Question 48:} D
\\ \textit{कमी होने से करेंसी का ज़्यादा बनना रुकता है, जिससे उसकी वैल्यू और खरीदने की ताकत समय के साथ बनी रहती है।}

\item \textbf{Question 49:} D
\\ \textit{पारंपरिक भुगतान प्रणालियाँ स्थानीय और अंतरराष्ट्रीय स्तर पर भुगतानकर्ताओं और प्राप्तकर्ताओं के बीच धन हस्तांतरण के लिए तंत्र और बुनियादी ढांचा प्रदान करती हैं।}

\item \textbf{Question 50:} C
\\ \textit{कागज़ के नोट, सिक्कों की तुलना में हल्के और आसानी से ले जाने में सुविधाजनक होने के कारण, लंबी दूरी के व्यापार को और अधिक कुशल बनाते हैं और ढुलाई की समस्याओं को कम करते हैं।}

\item \textbf{Question 51:} B
\\ \textit{नकदी के विपरीत, जिसका तुरंत आदान-प्रदान हो जाता है, क्रेडिट कार्ड लेनदेन बैंकों, कार्ड नेटवर्क और समाशोधन प्रक्रियाओं से जुड़े कई चरणों से गुजरता है, जिसके परिणामस्वरूप देरी होती है और अतिरिक्त लागत आती है।}

\item \textbf{Question 52:} A
\\ \textit{ज़्यादातर छोटे व्यवसायों को CSV एक्सपोर्ट की सरलता और पर्याप्त डिटेल से फायदा होता है।}

\item \textbf{Question 53:} D
\\ \textit{एक आम Bitcoin लेनदेन की CSV फाइल में आमतौर पर तारीख, रकम, फिएट समकक्ष, और Bitcoin पते का जिक्र होता है, लेकिन इस्तेमाल किए गए स्क्रिप्ट प्रकार का नहीं।}

\item \textbf{Question 54:} D
\\ \textit{इन ब्यौरों को दर्ज करने से स्पष्टता आती है और अकाउंटेंट को आने वाले और जाने वाले फंड्स में अंतर समझने में मदद मिलती है।}

\item \textbf{Question 55:} D
\\ \textit{स्पैरो Bitcoin अकाउंटिंग सॉल्यूशन नहीं है, जबकि बाकी सब हैं (यह एक Bitcoin Wallet सॉफ्टवेयर है)।}

\item \textbf{Question 56:} C
\\ \textit{ये टूल डेटा कलेक्शन, फॉर्मेट कन्वर्जन, टैक्स कैलकुलेशन और ट्रांजैक्शन कैटेगराइज़ेशन को ऑटोमेट करने के लिए बनाए गए हैं।}

\item \textbf{Question 57:} A
\\ \textit{नॉन-कस्टोडियल वॉलेट्स में प्राइवेट कीज़ और ट्रांजैक्शन का पूरा कंट्रोल यूज़र के हाथ में होता है।}

\item \textbf{Question 58:} A
\\ \textit{फिएट वैल्यू को ट्रैक करने से बिजनेस को Bitcoin की उतार-चढ़ाव वाली कीमतों के साथ बुककीपिंग मैनेज करने में मदद मिलती है।}

\item \textbf{Question 59:} C
\\ \textit{कस्टोडियल वॉलेट्स प्राइवेट कीज़ को मैनेज करते हैं, जिससे यूज़र्स के लिए टेक्निकल दिक्कतें कम हो जाती हैं।}

\item \textbf{Question 60:} A
\\ \textit{स्टार्टर प्रोफाइल कस्टोडियल या नॉन-कस्टोडियल लाइटनिंग वॉलेट्स का इस्तेमाल करती है।}

\item \textbf{Question 61:} B
\\ \textit{स्टार्टर प्रोफाइल के लिए सिर्फ एक स्मार्टफोन या टैबलेट चाहिए जिसमें लाइटनिंग Wallet हो, यह एक हल्का-फुल्का समाधान है।}

\item \textbf{Question 62:} A
\\ \textit{एसेंशियल प्रोफाइल में पीओएस ऐप कर्मचारियों को पहले से सेट मेनू से आइटम चुनने की सुविधा देता है, जिससे बिलिंग प्रक्रिया आसान हो जाती है।}

\item \textbf{Question 63:} B
\\ \textit{ओपन नोड को एक पूरी तरह से कस्टोडियल समाधान के रूप में पहचाना जाता है, जो स्विस जीडब्ल्यू-2 पे के हाइब्रिड स्वभाव से अलग है।}

\item \textbf{Question 64:} B
\\ \textit{एसेंशियल प्रोफाइल छोटे से मध्यम आकार के व्यवसायों के लिए है जिनके पास कम Bitcoin लेनदेन होते हैं, जैसे कि एक रेस्तरां।}

\item \textbf{Question 65:} C
\\ \textit{स्विस Bitcoin पे एक सीधे-सादे PoS Interface के साथ डिज़ाइन किया गया है जो रोज़मर्रा के कामों को आसान बनाता है।}

\item \textbf{Question 66:} A
\\ \textit{हम Swiss Bitcoin Pay को एक Essential प्रोफाइल के लिए सुझाया गया समाधान बताते हैं, जो सरलता और सुरक्षा का सही संतुलन बनाता है।}

\item \textbf{Question 67:} B
\\ \textit{बी-बॉप एक विकल्प है उन लोगों के लिए जो एक एकीकृत ऑनलाइन ई-स्टोर चाहते हैं जो जीडब्ल्यू-1 भुगतान की सुविधा देता है।}

\item \textbf{Question 68:} D
\\ \textit{ज़ैप्राइट और मस्केट चेकआउट प्रक्रिया में कस्टमाइज़ेशन और रिपोर्टिंग सुधार जोड़ते हैं, जो मुख्य भुगतान बुनियादी ढांचे को पूरक बनाते हैं।}

\item \textbf{Question 69:} B
\\ \textit{BTC Pay Server को ऑन-प्रिमाइसेस या क्लाउड पर डिप्लॉय करने की लचीलापन व्यवसायों को उनकी परिचालन आवश्यकताओं के अनुरूप सबसे उपयुक्त सेटअप चुनने की सुविधा देता है।}

\item \textbf{Question 70:} C
\\ \textit{कोर्स में BTC पे सर्वर को प्रोफेशनल सेटअप्स के लिए कोर ओपन-सोर्स सॉल्यूशन के रूप में हाइलाइट किया गया है, क्योंकि इसमें एक्सटेंसिव इंटीग्रेशन ऑप्शंस और कंट्रोल होते हैं।}

\item \textbf{Question 71:} D
\\ \textit{प्रोफेशनल बिज़नेस को अपने जटिल वर्कफ़्लो के हिसाब से कस्टमाइज़ेबल इनवॉइस, सोफिस्टिकेटेड रिपोर्टिंग, और रोल मैनेजमेंट की ज़रूरत होती है।}

\item \textbf{Question 72:} D
\\ \textit{एलएसपी (LSPs) लाइटनिंग पर जटिल और बड़े पैमाने के ऑपरेशन्स को संभालते हैं।}

\item \textbf{Question 73:} A
\\ \textit{वर्कफ्लो ट्रिगर्स सिस्टम को रियल टाइम में अपडेट करते हैं।}

\item \textbf{Question 74:} A
\\ \textit{APIs से लगातार, रियल-टाइम इंटीग्रेशन संभव होता है।}

\item \textbf{Question 75:} B
\\ \textit{बड़ी रकम को मैनेज करने के लिए मल्टी-सिग्नेचर कॉन्फिगरेशन और टाइमलॉक जैसी हाई-सिक्योरिटी वाली सुविधाएं बेहद जरूरी हैं।}

\item \textbf{Question 76:} A
\\ \textit{एंटरप्राइज़ सॉल्यूशन्स का मकसद है Bitcoin पेमेंट्स को कोर बिज़नेस वर्कफ़्लोज़ और सिस्टम्स के साथ जोड़ना।}


\end{enumerate}

\end{document}
